\documentclass[12pt,a4paper,reqno]{amsart}

% language
\usepackage[greek,english]{babel}
\usepackage[utf8]{inputenc}
\usepackage{alphabeta}

% change default names to greek
\addto\captionsenglish{
    \renewcommand{\contentsname}{Περιεχόμενα}
    \renewcommand{\refname}{Βιβλιογραφία}
    \renewcommand{\datename}{Ημερομηνία:}
    \renewcommand{\urladdrname}{Ιστοσελίδα}
}

% math 
\usepackage{amsmath,amsthm,amssymb,amscd}

% font
\usepackage[cal=euler]{mathalfa}
\usepackage{libertinus-type1}
% \usepackage{txfonts} % for upright greek letters
\usepackage{bm} % for bold symbols
\usepackage{bbm} % for the simply-looking bb symbols

% miscellaneous 
\usepackage{changepage} %for indenting environments
\usepackage{csquotes} % example: \textcquote{}
\usepackage{blkarray}
\setcounter{MaxMatrixCols}{20} % default for pmatrix is 10!!
\usepackage{ytableau}
\usepackage{array} %needed to increase the vertical length in a tabular

% drawing
\usepackage{tikz,tikz-cd}
\usetikzlibrary{shapes.misc, patterns, matrix, calc, intersections,positioning,arrows.meta}
\usepackage{graphics,graphicx}
\usepackage{float} % provides enhanced control and customization options for floating objects such as figures and tables

% colors
\usepackage{xcolor}
\definecolor{darkcandyapplered}{rgb}{0.64, 0.0, 0.0}
\definecolor{midnightblue}{rgb}{0.1, 0.1, 0.44}
\definecolor{mylightblue}{HTML}{336699}
\definecolor{burntorange}{rgb}{0.8, 0.33, 0.0}
\definecolor{iceberg}{rgb}{0.44, 0.65, 0.82}
\definecolor{applegreen}{rgb}{0.55, 0.71, 0.0}
\definecolor{canaryyellow}{rgb}{1.0, 0.94, 0.0}
\definecolor{lightgray}{rgb}{0.83, 0.83, 0.83}

% hrefs
\usepackage{hyperref}
\usepackage[noabbrev,capitalize]{cleveref}
\hypersetup{
    pdftoolbar=true,        
    pdfmenubar=true,        
    pdffitwindow=false,     
    pdfstartview={FitH},  % fits the width of the page to the window
    pdftitle={},
    pdfauthor={},
    pdfsubject={},
    pdfkeywords={},
    pdfnewwindow=true,  % links in new window
    colorlinks=true,  % false: boxed links; true: colored links
    linkcolor=darkcandyapplered,   % color of internal links
    citecolor=midnightblue,  % color of links to bibliography
    urlcolor=cyan,  % color of external links
    linktocpage=true  % changes the links from the section body to the page number
    }

% geometry
\textwidth=16cm 
\textheight=21cm 
\hoffset=-55pt 
\footskip=25pt

% thm envs (you might need to change the path)
\theoremstyle{definition}
\newtheorem{exercise}{Άσκηση}
\newtheorem{solution}{Λύση}
\newtheorem*{remark}{Παρατήρηση}
\newtheorem*{example}{Παράδειγμα}

% math ops 
% fields
\newcommand{\NN}{\mathbbmss{N}} 
\newcommand{\ZZ}{\mathbbmss{Z}} 
\newcommand{\QQ}{\mathbbmss{Q}} 
\newcommand{\RR}{\mathbbmss{R}} 
\newcommand{\CC}{\mathbbmss{C}} 
\newcommand{\KK}{\mathbbmss{K}} 
\newcommand{\FF}{\mathbbmss{F}} 
\newcommand{\HH}{\mathbbmss{H}} 

% symmetric group
\newcommand{\fS}{\mathfrak{S}}  

% calligraphic 
\newcommand{\aA}{\mathcal{A}} 
\newcommand{\bB}{\mathcal{B}}
\newcommand{\cC}{\mathcal{C}}
\newcommand{\dD}{\mathcal{D}}
\newcommand{\eE}{\mathcal{E}}
\newcommand{\fF}{\mathcal{F}}
\newcommand{\hH}{\mathcal{H}}
\newcommand{\iI}{\mathcal{I}}
\newcommand{\lL}{\mathcal{L}}
\newcommand{\oO}{\mathcal{O}}
\newcommand{\pP}{\mathcal{P}}
\newcommand{\sS}{\mathcal{S}}
\newcommand{\mM}{\mathcal{M}}
\newcommand{\uU}{\mathcal{U}}

% bold
\newcommand{\bfa}{\mathbf{a}}
\newcommand{\bfe}{\mathbf{e}}
\newcommand{\bfF}{\pmb{F}}
\newcommand{\bfR}{\pmb{R}}
\newcommand{\bfv}{\mathbf{v}}
%\newcommand{\bfx}{\bm{x}}
%\newcommand{\bfx}{\mathbf{x}} 
\newcommand{\bfx}{\pmb{x}}
\newcommand{\bfX}{\pmb{X}}
\newcommand{\bfy}{\pmb{y}}
\newcommand{\bfz}{\pmb{z}}

% roman
\newcommand{\rmA}{\mathrm{A}}
\newcommand{\rmB}{\mathrm{B}}
\newcommand{\rmC}{\mathrm{C}}
\newcommand{\rmD}{\mathrm{D}} 
\newcommand{\rmH}{\mathrm{H}} 
\newcommand{\rmh}{\mathrm{h}} 
\newcommand{\rmI}{\mathrm{I}} 
\newcommand{\rmK}{\mathrm{K}}
\newcommand{\rmM}{\mathrm{M}}
\newcommand{\rmP}{\mathrm{P}}  
\newcommand{\rmp}{\mathrm{p}}  
\newcommand{\rmQ}{\mathrm{Q}}  
\newcommand{\rmR}{\mathrm{R}}
\newcommand{\rmS}{\mathrm{S}}
\newcommand{\rmT}{\mathrm{T}}
\newcommand{\rmU}{\mathrm{U}}
\newcommand{\rmV}{\mathrm{V}}
\newcommand{\rmY}{\mathrm{Y}}
\newcommand{\rmZ}{\mathrm{Z}}
\newcommand{\rmz}{\mathrm{z}}

% arrows and symbols 
\renewcommand{\to}{\rightarrow}
\newcommand{\toto}{\longrightarrow}
\newcommand{\mapstoto}{\longmapsto}
\newcommand{\then}{\Rightarrow}
\newcommand{\IFF}{\Leftrightarrow}
\newcommand{\tl}{\tilde}
\newcommand{\wtl}{\widetilde}
\newcommand{\ol}{\overline}
\newcommand{\ul}{\underline}
\newcommand{\oldemptyset}{\emptyset}
\renewcommand{\emptyset}{\varnothing}
\DeclareMathSymbol{\Arg}{\mathbin}{AMSa}{"39} % for arguments 
\newcommand{\onto}{\ensuremath{\twoheadrightarrow}}
\newcommand{\tle}{\trianglelefteq}
\newcommand{\tge}{\trianglerighteq}

% custom ops
\newcommand{\rmKer}{\mathrm{Ker}}
\newcommand{\rmIm}{\mathrm{Im}}
\newcommand{\lcm}{\mathrm{lcm}}
\newcommand{\GL}{\mathrm{GL}}
\newcommand{\ev}{\mathrm{ev}}
\newcommand{\End}{\mathrm{End}}
\newcommand{\rmid}{\mathrm{id}}
\newcommand{\rmspan}{\mathrm{span}}

% absolute value symbol
\usepackage{mathtools} 
\DeclarePairedDelimiter\abs{\lvert}{\rvert}%
\DeclarePairedDelimiter\norm{\lVert}{\rVert}%
\makeatletter
\let\oldabs\abs
\def\abs{\@ifstar{\oldabs}{\oldabs*}}

% tensor symbol
\newcommand{\tensor}[1]{%
  \mathbin{\mathop{\otimes}\limits_{#1}}%
}

% permutation cycle notation
\ExplSyntaxOn
\NewDocumentCommand{\cycle}{ O{\;} m }
 {
  (
  \alec_cycle:nn { #1 } { #2 }
  )
 }

\seq_new:N \l_alec_cycle_seq
\cs_new_protected:Npn \alec_cycle:nn #1 #2
 {
  \seq_set_split:Nnn \l_alec_cycle_seq { , } { #2 }
  \seq_use:Nn \l_alec_cycle_seq { #1 }
 }
\ExplSyntaxOff

% setminus symbol
\newcommand{\mysetminusD}{\hbox{\tikz{\draw[line width=0.6pt,line cap=round] (3pt,0) -- (0,6pt);}}}
\newcommand{\mysetminusT}{\mysetminusD}
\newcommand{\mysetminusS}{\hbox{\tikz{\draw[line width=0.45pt,line cap=round] (2pt,0) -- (0,4pt);}}}
\newcommand{\mysetminusSS}{\hbox{\tikz{\draw[line width=0.4pt,line cap=round] (1.5pt,0) -- (0,3pt);}}}
\newcommand{\sm}{\mathbin{\mathchoice{\mysetminusD}{\mysetminusT}{\mysetminusS}{\mysetminusSS}}}

% miscellaneous commands
\newcommand{\defn}[1]{{\color{mylightblue}{#1}}}
\newcommand{\toDo}{{\bf\color{red} TODO}}
\newcommand{\toCite}{{\bf\color{green} CITE}}
\newcommand*{\vertbar}{\rule[-1ex]{0.5pt}{2.5ex}} % for matrices with column vectors
\newcommand*{\horzbar}{\rule[.5ex]{2.5ex}{0.5pt}} % for matrices with row vectors
\newcommand{\myblue}[1]{{\color{iceberg}{#1}}}
\newcommand{\myorange}[1]{{\color{burntorange}{#1}}}
\newcommand{\mygreen}[1]{{\color{applegreen}{#1}}}
\newcommand{\myred}[1]{{\color{darkcandyapplered}{#1}}}

% ferrer's diagram
\newcommand{\fdiagram}[1]{
    \begin{tikzpicture}[scale=.7]
        \fill foreach \Z [count=\Y] in {#1}
        {foreach \X in {1,...,\Z} 
        {(\X,-\Y) circle[radius=3pt]}};
    \end{tikzpicture}
}

%
\usepackage{pgfplots}
\pgfplotsset{compat=1.18}

%
\makeatletter
\def\mathcolor#1#{\@mathcolor{#1}}
\def\@mathcolor#1#2#3{%
  \protect\leavevmode
  \begingroup
    \color#1{#2}#3%
  \endgroup
}
\makeatother

% 
\newenvironment{nouppercase}{%
  \let\uppercase\relax%
  \renewcommand{\uppercasenonmath}[1]{}}{}

% titlepage
\title{226: Θεωρία Δακτυλίων και Modules \\ Ασκήσεις με Ενδεικτικές Λύσεις}
\author[Β.~Δ. Μουστακας]{Βασίλης Διονύσης Μουστάκας \\ Πανεπιστήμιο Κρήτης}
\date{25 Φεβρουαρίου 2026}
% \urladdr{\href{https://sites.google.com/view/vasmous}{https://sites.google.com/view/vasmous}}

\begin{document}

\begingroup
\def\uppercasenonmath#1{} % this disables uppercase title
\let\MakeUppercase\relax % this disables uppercase authors
\maketitle
\endgroup

% \setcounter{section}{1}
% \setcounter{theorem}{5}
% \setcounter{equation}{}
\renewcommand{\thesubsection}{\thesection.\arabic{subsection}}

Σε ότι ακολουθεί υποθέτουμε ότι $R$ ένας δακτύλιος και $n$ είναι ένας θετικός ακέραιος, εκτός αν αναφέρεται διαφορετικά.

\begin{exercise}{\rm(\cite[Ασκήσεις~7.1.7~και~7.2.7]{DF04})}
\label{ex:DF_7.1.7}
Το σύνολο
\[
\rmZ(R) \coloneqq \{a \in R : ra=ar, \ \text{για κάθε $r \in R$}\}
\]
ονομάζεται \defn{κέντρο} του $R$. Να δείξετε τα εξής.
\begin{itemize}
  \item[(1)] Το $\rmZ(R)$ είναι υποδακτύλιος του $R$.
  \item[(2)] Αν ο $R$ είναι δακτύλιος διαίρεσης, τότε το $\rmZ(R)$ είναι σώμα.
  \item[(3)] Έχουμε $\rmZ(\rmM_n(R)) = \{a \, I_n : a \in \rmZ(R)\}$, όπου $I_n$ είναι ο ταυτοτικός $n\times{n}$-πίνακας.
\end{itemize}
\end{exercise}  

\begin{proof}[Λύση]
  \leavevmode
  \begin{itemize}
    \item[(1)] Θα επαληθεύσουμε τον Ορισμό 1.3. Το κέντρο περιέχει το μηδέν και την μονάδα του $R$, καθώς 
    \begin{align*}
    r0 &= 0r = 0 \\
    r1 &= 1r = r 
    \end{align*}
    για κάθε $r \in R$. Επιπλέον, είναι κλειστό ως προς την πρόσθεση και ως προς αντίθετα στοιχεία, διότι 
    \[
    (a-b)r = ar - br = ra - rb = r(a-b),
    \]
    για κάθε $a, b \in \rmZ(R)$ και $r \in R$. Τέλος, είναι κλειστό ως προς τον πολλαπλασιασμό, διότι 
    \[
    abr = arb = rab
    \]
    για κάθε $a, b \in \rmZ(R)$ και $r \in R$.
    \item[(2)] Το $\rmZ(R)$ είναι προφανώς μεταθετικός υποδακτύλιος του $R$. Γι αυτό, αρκεί να δείξουμε ότι κάθε μη μηδενικό στοιχείο του έχει αντίστροφο ως προς τον πολλαπλασιασμό. Έστω $a \in \rmZ(R) \subseteq R$. Επειδή ο $R$ είναι δακτύλιος διαίρεσης (βλ. Ορισμό 1.1) υπάρχει το $a^{-1} \in R$. Παρατηρούμε ότι 
    \[
    ra^{-1} = (a^{-1}a)ra^{-1} = a^{-1}(ar)a^{-1} = a^{-1}r(aa^{-1}) = a^{-1}r
    \]
    για κάθε $r \in R$ και γι αυτό $a^{-1} \in \rmZ(R)$ και το ζητούμενο έπεται.
    \item[(γ)] Για κάθε $1 \le i, j \le n$, θεωρούμε τον πίνακα $E_{ij} \in \rmM_n(R)$ ο οποίος έχει $1$ στη θέση $(i,j)$ και $0$ παντού αλλού. Για κάθε $A \in \rmM_n(R)$, παρατηρούμε ότι 
    \begin{itemize}
      \item[$\bullet$] ο $AE_{ij}$ είναι ο πίνακας του οποίου η $j$-οστή στήλη ισούται με την $i$-οστή στήλη του $A$ και όλες οι υπόλοιπες είναι $0$, και 
      \item[$\bullet$] ο $E_{ij}A$ είναι ο πίνακας του οποίου η $i$-οστή γραμμή ισούται με την $j$-οστή γραμμή του $A$ και όλες οι υπόλοιπες είναι $0$
    \end{itemize}
    (γιατί;). Για παράδειγμα, για $n=4$ και $(i,j) = (2,3)$ 
    \begin{align*}
      \begin{pmatrix}
        a_{11} & a_{12} & a_{13} & a_{14} \\
        a_{21} & a_{22} & a_{23} & a_{24} \\
        a_{31} & a_{32} & a_{33} & a_{34} \\
        a_{41} & a_{42} & a_{43} & a_{44} 
      \end{pmatrix}
      \begin{pmatrix}
        0 & 0 & 0 & 0 \\
        0 & 0 & 1 & 0 \\
        0 & 0 & 0 & 0 \\
        0 & 0 & 0 & 0 
      \end{pmatrix}
      &=
      \begin{pmatrix}
        0 & 0 & a_{12} & 0 \\
        0 & 0 & a_{22} & 0 \\
        0 & 0 & a_{32} & 0 \\
        0 & 0 & a_{42} & 0 
      \end{pmatrix} \\
      \begin{pmatrix}
        0 & 0 & 0 & 0 \\
        0 & 0 & 1 & 0 \\
        0 & 0 & 0 & 0 \\
        0 & 0 & 0 & 0 
      \end{pmatrix}
      \begin{pmatrix}
        a_{11} & a_{12} & a_{13} & a_{14} \\
        a_{21} & a_{22} & a_{23} & a_{24} \\
        a_{31} & a_{32} & a_{33} & a_{34} \\
        a_{41} & a_{42} & a_{43} & a_{44} 
      \end{pmatrix}
      &=
      \begin{pmatrix}
        0 & 0 & 0 & 0 \\
        0 & 0 & 0 & 0 \\
        a_{31} & a_{32} & a_{33} & a_{34} \\
        0 & 0 & 0 & 0
      \end{pmatrix}.
    \end{align*} 
    
    Συνεπώς, αν $A \in \rmZ(\rmM_n(R))$, τότε $AE_{ij} = E_{ij}A$, για κάθε $1 \le i,j \le n$, που σημαίνει ότι έχει την μορφή 
    \[
    A = a \, I_n
    \]
    για κάποιο $a \in R$ (γιατί;).  Επιπλέον, για κάθε $r \in R$, 
    \[
    (ra)I_n = (r\,I_n)A = A(r\,I_n) = (ar)I_n,
    \]
    όπου η δεύτερη ισότητα έπεται από το ότι $A \in \rmZ(\rmM_n(R))$ και γι αυτό $ra = ar$ ή ισοδύναμα $a \in \rmZ(R)$. Επομένως,
    \[
    \rmZ(\rmM_n(R)) \subseteq \{a \, I_n : a \in \rmZ(R)\}.
    \]
    Ο αντίστροφος εγκλεισμός είναι προφανής και το ζητούμενο έπεται.
  \end{itemize}
\end{proof}

\begin{remark}
  Στην περίπτωση όπου ο $R$ είναι δακτύλιος διαίρεσης και το κέντρο του είναι σώμα, έστω $F$, ο $R$ γίνεται διανυσματικός χώρος πάνω από το $F$ με βαθμωτό πολλαπλασιασμό
  \begin{align*}
    F\times R &\to R \\
        (a,r) &\mapsto ar.
  \end{align*}
  Επιπλέον, 
  \begin{equation}
    \label{eq:algebra}
    a\cdot(rs) = (a\cdot{r})s = r(a\cdot{s}) 
  \end{equation}
  για κάθε $a \in F$ και $r, s \in R$. Ένας δακτύλιος που έχει ταυτόχρονα δομή διανυσματικού χώρου με τους δύο πολλαπλασιασμούς να είναι συμβατοί με την έννοια της Ταυτότητας \eqref{eq:algebra} ονομάζεται \defn{$F$-άλγεβρα}. Για παράδειγμα, ο πολυωνυμικός δακτύλιος $\RR[x]$ είναι μια $\RR$-άλγεβρα. 
\end{remark}

\begin{exercise}{\rm(βλ. \cite[Άσκηση~7.2.12]{DF04})}
  \label{ex:groupRing}
  Έστω $F$ ένα σώμα και $G = \{g_1, g_2, \dots, g_n\}$ μία πεπερασμένη ομάδα με ταυτοτικό στοιχείο $\epsilon$. Έστω $FG$ ο διανυσματικός χώρος πάνω από το $F$ με βάση τα στοιχεία της $G$. Τα στοιχεία του $FG$ είναι της μορφής\footnote{Αυτό ονομάζεται \emph{τυπικό άθροισμα} των $g_1, g_2, \dots, g_n$.} 
  \[
  \lambda_1g_1 + \lambda_2g_2 + \cdots + \lambda_ng_n 
  \]
  για κάποια $\lambda_1, \lambda_2, \dots, \lambda_n \in F$. Το μηδέν του $FG$ είναι το στοιχείο
  \[
  0_{FG} \coloneqq 0_Fg_1 + 0_Fg_2 + \cdots + 0_Fg_n. 
  \]
  Στο $FG$ ορίζουμε έναν πολλαπλασιασμό θέτοντας 
  \[
  (\lambda g_i) \cdot (\mu g_j) = (\lambda\mu) g_k
  \]
  για κάθε $\lambda, \mu \in F$ και $1 \le i, j \le n$, όπου $g_k = g_ig_j \in G$ και επεκτείνοντας γραμμικά σε ολόκληρο το $FG$. Να δείξετε τα εξής.
  \begin{itemize}
    \item[(1)] Ο $FG$ είναι δακτύλιος με μονάδα $1_{FG} \coloneqq 1_F\epsilon$, ο οποίος ονομάζεται \defn{ομαδοδακτύλιος} (group ring) της $G$.
    \item[(2)] Ο $FG$ είναι μεταθετικός δακτύλιος αν και μόνο αν η $G$ είναι αβελιανή ομάδα.
    \item[(2)] Αν $\abs{G} > 1$, τότε ο $FG$ έχει μηδενοδιαιρέτες.
    \item[(3)] Έχουμε $g_1 + g_2 + \cdots + g_n \in \rmZ(FG)$. 
  \end{itemize}
\end{exercise}

\begin{example}
Για παράδειγμα, έστω
\[
G = \rmS_3 = \{ \epsilon, \cycle{1,2}, \cycle{2,3}, \cycle{1,3}, \cycle{1,2,3}, \cycle{1,3,2}\}
\]
η συμμετρική ομάδα τάξης $6$ και $F = \RR$. Ο $\RR\rmS_3$ ως διανυσματικός χώρος πάνω από το $\RR$ έχει διάσταση $6$. θεωρούμε τα ακόλουθα στοιχεία του $\RR\rmS_3$
\begin{align*}
  v &= 3\cycle{1,2} - 5\cycle{2,3} + 14\cycle{1,2,3} \\
  u &= 6\epsilon + 2\cycle{2,3} + 7\cycle{1,3,2}.
\end{align*}
Υπολογίζουμε (βλ. \cite[Άσκηση~7.2.10]{DF04})
\begin{align*}
  v + u 
  &= 6\epsilon + 3\cycle{1,2} - 3\cycle{2,3} + 14\cycle{1,2,3} + 7\cycle{1,3,2} \\
  2v - 3u 
  &= 6\cycle{1,2} - 10\cycle{2,3} + 28\cycle{1,2,3} - 18\epsilon + -6\cycle{2,3} - 21\cycle{1,3,2} \\
  &= -18\epsilon + 6\cycle{1,2} -16\cycle{2,3} + 28\cycle{1,2,3} - 21\cycle{1,3,2} \\
  vu &= \left(
    3\cycle{1,2} - 5\cycle{2,3} + 14\cycle{1,2,3}
  \right)
  \left(
    6\epsilon + 2\cycle{2,3} + 7\cycle{1,3,2}
  \right) \\
  &= 18\cycle{1,2} + 63\cycle{1,2}\cycle{2,3} + 21\cycle{2,3}\cycle{1,3,2} -30\cycle{2,3} -10\cycle{2,3}^2 -35\cycle{2,3}\cycle{1,3,2} \\
  &\quad + 84\cycle{1,2,3} + 28\cycle{1,2,3}\cycle{2,3} + 98\cycle{1,2,3}\cycle{1,3,2} \\
  &= 18\cycle{1,2} + 63\cycle{1,2,3} + 21\cycle{1,2} -30\cycle{2,3} -10\epsilon -35\cycle{1,2} \\
  &\quad + 84\cycle{1,2,3} + 28\cycle{1,2} + 98\epsilon \\
  &= 88\epsilon + 32\cycle{1,2} - 30\cycle{2,3} + 147\cycle{1,2,3}.
\end{align*}
Ποιά είναι τα $uv$ και $v^2$;
\end{example}

\begin{proof}[Λύση της Άσκησης \ref{ex:groupRing}]
  Τα (1) και (2) αφήνονται ως άσκηση.
  \begin{itemize}
    \item[(2)] Αν $\abs{G} > 1$, τότε υπάρχει κάποιο $g \in G$ με $g \neq \epsilon$. Αν $m$ είναι η τάξη του $g$, τότε 
    \[
    (\underbrace{\epsilon - g}_{\neq 0_{FG}})(\underbrace{\epsilon + g + \cdots + g^{m-1}}_{\neq 0_{FG}}) = \epsilon - g^m = \epsilon - \epsilon = 0_{FG}.
    \]
    Συνεπώς, το $\epsilon - g$ είναι μηδενοδιαιρέτης του $FG$.
    \item[(3)] Έστω $a = g_1 + g_2 + \cdots + g_n$. Αρκεί να δείξουμε ότι $ag_i = g_ia$, για κάθε $1 \le i \le n$ (γιατί;). Πράγματι,
    \[
    (g_1 + g_2 + \cdots + g_n)g_i = g_1 + g_2 + \cdots + g_n,
    \]
    διότι η απεικόνιση 
    \begin{align*}
      G &\to G \\
      g &\mapsto gg_i
    \end{align*}
    είναι 1-1 και επί, για κάθε $1 \le i \le n$. Ομοίως, 
    \[
    g_i(g_1 + g_2 + \cdots + g_n) = g_1 + g_2 + \cdots + g_n,
    \]
    για κάθε $1\le i \le n$ και το ζητούμενο έπεται.
  \end{itemize}
\end{proof}

\begin{remark}
  Ο ομαδοδακτύλιος είναι το βασικό αντικείμενο μελέτης της \emph{θεωρίας αναπαραστάσεων}.
\end{remark}

\begin{exercise}{\rm(\cite[Ασκήσεις~7.3.25]{DF04})}
  \label{ex:7.3.25}
  Υποθέτουμε ότι ο $R$ είναι μεταθετικός δακτύλιος. Να δείξετε ότι 
  \[
  (a+b)^n = \sum_{k=0}^n \binom{n}{k} a^kb^{n-k},
  \]
  για κάθε $a, b \in R$ και κάθε $n \ge 0$, όπου με 
  \begin{equation} 
    \label{eq:binomial_thm}
    \binom{n}{k} \coloneqq \frac{n!}{k!(n-k)!}
  \end{equation}
  συμβολίζουμε τον \emph{διωνυμικό συντελεστή}.
\end{exercise}

\begin{proof}[Λύση]
  Έστω $a, b \in R$. Θα κάνουμε επαγωγή ως προς $n$. Για $n=0$ και τα δύο μέλη εξ' ορισμού είναι ίσα με $1_R$. Αυτή είναι η βάση της επαγωγής. Έστω $n > 0$ και υποθέτουμε ότι η Ταυτότητα \eqref{eq:binomial_thm} ισχύει για $n$. Αυτή είναι η επαγωγική υπόθεση. Τότε 
  \begin{align*}
    (a+b)^{n+1} &= (a+b)^n(a+b) \\ 
    &= \left(
      \sum_{k=0}^n \binom{n}{k} a^kb^{n-k}
    \right)(a+b) \\
    &= 
    \sum_{k=0}^n \binom{n}{k} a^{k+1}b^{n-k} + \sum_{k=0}^n \binom{n}{k} a^kb^{n-k+1} \\
    &=
    b^{n+1} + \sum_{k=1}^{n+1} \left(\binom{n}{k-1} + \binom{n}{k}\right) a^kb^{n+1-k} + a^{n+1} \\
    &= \sum_{k=0}^{n+1} \binom{n+1}{k} a^kb^{n+1-k},
  \end{align*}
  όπου η δεύτερη ισότητα έπεται από την επαγωγική υπόθεση και η τελευταία ισότητα έπεται από την \emph{ταυτότητα του Pascal}.
\end{proof}

\begin{remark}
  Η Άσκηση \ref{ex:7.3.25} ισχύει για μη μεταθετικούς δακτύλιους με την ασθενέστερη υπόθεση ότι $ab = ba$.
\end{remark}

\begin{exercise}{\rm(\cite[Ασκήσεις~7.1.13]{DF04})}
  Ένα στοιχείο $x \in R$ ονομάζεται \defn{μηδενοδύναμο} (nilpotent) αν $x^m=0$ για κάποιο $m \ge 1$. Να δείξετε τα εξής.
  \begin{itemize}
    \item[(1)] Αν $n = a^kb$ για κάποιους θετικούς ακεραίους $a, b, k$, τότε το $ab$ είναι μηδενοδύναμο στο $\ZZ_n$.
    \item[(2)] Έστω $x \in \ZZ$. Το $x$ είναι μηδενοδύναμο στο $\ZZ_n$ αν και μόνο αν κάθε πρώτος διαιρέτης του $n$ διαιρεί το $x$.
    \item[(3)] Ποιά είναι τα μηδενοδύναμα στοιχεία του $\ZZ_{72}$;
  \end{itemize}
\end{exercise}

\begin{proof}[Λύση]
  \leavevmode
  \begin{itemize}
    \item[(1)] Έχουμε 
    \[
    (ab)^k = a^kb^k = \underbrace{a^kb}_{=n}b^{k-1} \equiv 0 \pmod{n}
    \]
    και γι αυτό το $ab$ είναι μηδενοδύναμο στο $\ZZ_n$.
    \item[(2)] Έστω $x \in \ZZ_n$ ένα μηδενοδύναμο στοιχείο και $p$ ένας πρώτος διαιρέτης του $n$. Τότε $x^m \equiv 0 \pmod{n}$ για κάποιο $m \ge 1$, ή ισοδύναμα $n \mid x^m$. Συνεπώς, $p \mid x^m$ και γι αυτό $p \mid x$, από το \emph{λήμμα του Ευκλείδη}\footnote{Το λήμμα του Ευκλείδη μας πληροφορεί ότι αν $p$ είναι πρώτος που διαιρεί το γινόμενο $ab$ δύο ακεραίων, τότε $p \mid a$ ή $p \mid b$.}.
    
    Αντιστρόφως, υποθέτουμε ότι κάθε πρώτος διαιρέτης του $n$ διαιρεί το $x \in \ZZ$. Συνεπώς, αν $n = p_1^{k_1}p_2^{k_2}\cdots p_\ell^{k_\ell}$ είναι η παραγοντοποίηση του $n$ σε πρώτους, τότε 
    \[
    x = (p_1p_2\cdots p_\ell)b
    \]
    για κάποιο $b \ge 1$. Διαλέγοντας $m = \max\{k_1, k_2, \dots, k_\ell\}$ έχουμε 
    \[
    x^m = (p_1p_2\cdots p_\ell)^mb^m = n cb^m \equiv 0 \pmod{n}
    \]
    για κάποιο $c \ge 1$. Οπότε, το $x$ είναι μηδενοδύναμο στο $\ZZ_n$ και το ζητούμενο έπεται.
    \item[(3)] Έχουμε $72 = 2^33^2$ και γι αυτό, από το (2), τα μηδενοδύναμα στοιχεία του $\ZZ_{72}$ είναι τα κοινά πολλαπλάσια του $2$ \emph{και} του $3$, δηλαδή
    \[
    0, 6,12,18,24,30,36,42,48,54,60,66.
    \]
  \end{itemize}
\end{proof}

\begin{thebibliography}{99}
    \bibitem{DF04}
    D. S. Dummit και R. M. Foote, 
    \textit{Abstract algebra},
    third edition, Wiley, 2004. 
\end{thebibliography}
\end{document}